\chapter{Introduction}

\section{Background and Motivation}

Serverless computing has emerged as a dominant cloud paradigm, abstracting infrastructure management and enabling developers to focus on application logic \citep{Shafiei2022Survey, Wen2023RiseServerless}. Function-as-a-Service (FaaS) platforms such as AWS Lambda execute stateless event-driven functions on ephemeral compute instances, scaling automatically in response to demand. However, distributed systems are inherently subject to transient faults---network partitions, timeouts, and service unavailability---prompting serverless platforms to adopt \textit{at-least-once} invocation semantics to ensure progress \citep{AWSLambdaDocs}.

At-least-once invocation guarantees that a function executes successfully at least once per event, but platform retries introduce the risk of \textit{duplicate state mutations}: if a function modifies external state (e.g., writes to a database) and subsequently times out or fails after the state change has committed, a retry will attempt the operation again, potentially yielding inconsistent or corrupted data. The burden of idempotency---ensuring that duplicate executions produce the same result as a single execution---is delegated to the application developer \citep{Shafiei2022Survey, Wen2023RiseServerless}.

\section{The Landscape of Solutions}

Recent research has proposed custom serverless runtimes to eliminate duplicate mutations entirely. \citet{Qi2025Halfmoon} present Halfmoon, which employs asymmetric logging on a custom runtime to achieve exactly-once state mutation semantics with minimal logging overhead. Similarly, Boki \citep{Jia2024Boki}, CausalMesh \citep{Zhang2024CausalMesh}, LambdaStore \citep{Mast2026LambdaStore}, Styx \citep{Psarakis2025Styx}, Netherite \citep{Burckhardt2022Netherite}, and $\mu$2sls \citep{Kallas2023Mu2sls} provide strong consistency guarantees by replacing or augmenting the underlying FaaS runtime. While these systems demonstrate significant theoretical and empirical advances, they require platform modification or custom deployments---options unavailable to practitioners on managed services such as AWS Lambda, Google Cloud Functions, or Azure Functions.

For developers constrained to managed platforms, serverless providers expose \textit{application-level primitives} to mitigate duplicate mutations. Amazon DynamoDB, for instance, supports \textit{conditional writes} (PutItem or UpdateItem with ConditionExpression clauses) that atomically check preconditions before applying mutations \citep{Elhemali2022DynamoDB, AWSDynamoDBDocs}. AWS documentation further recommends the \textit{idempotency-key pattern}, wherein a function persists a unique request identifier and response payload, returning the cached result on subsequent deliveries of the same request \citep{AWSPowertoolsDocs}.

Despite widespread industry adoption of these patterns, \textbf{no prior empirical study has measured the duplicate-mutation reduction, end-to-end latency overhead, and consumed capacity costs of conditional writes and idempotency-key patterns on stock AWS Lambda and DynamoDB under a known, controlled retry schedule}. Existing literature either focuses on custom runtime correctness \citep{Qi2025Halfmoon, Jia2024Boki, Zhang2024CausalMesh} or formal verification of idempotence properties \citep{Ding2023Flux} without quantifying the practical trade-offs on managed infrastructure.

\section{Research Question and Objectives}

This research addresses the following question:

\begin{quote}
\textit{By how much do a conditional write (P2) and an idempotency-key write (P3) reduce duplicate state mutation, relative to a plain unconditional write (P1), when an AWS Lambda function is retried under injected timeouts---and what do they cost in end-to-end latency and consumed DynamoDB capacity?}
\end{quote}

\noindent The study pursues four objectives:

\begin{enumerate}
    \item Quantify duplicate-mutation rates for P1 (plain put), P2 (conditional put with \texttt{attribute\_not\_exists} or version guard), and P3 (idempotency-key pattern) under known injected retry multiplicities of 1, 2, and 5 deliveries per request.
    \item Measure mean and 95th-percentile end-to-end latency for each write path, distinguishing warm invocations from cold starts.
    \item Measure consumed DynamoDB capacity (read and write units) per request, including failed conditional checks (which are billed despite not mutating state).
    \item Position each write path on the \textit{correctness--cost surface}---plotting duplicate-mutation rate against latency and capacity overhead---and contrast the findings with the baseline idea from \citet{Qi2025Halfmoon} (no duplicate mutation after retry via asymmetric logging on a custom runtime).
\end{enumerate}

\section{Significance and Beneficiaries}

The findings are of practical interest to three constituencies:

\begin{itemize}
    \item \textbf{Cloud application developers} constrained to managed FaaS platforms, who must choose among application-level idempotency strategies without empirical cost--benefit data.
    \item \textbf{Serverless platform architects}, who design managed services and set default retry policies, benefit from understanding the trade-offs users face when delegating idempotency to the application layer.
    \item \textbf{Researchers} in distributed systems and cloud computing, who can use the measured correctness--cost surface to contextualise custom-runtime proposals against the baseline capabilities of existing managed platforms.
\end{itemize}

\section{Scope and Limitations}

This study evaluates \textit{single-item writes} on \textit{Amazon DynamoDB on-demand capacity mode} from \textit{AWS Lambda Python 3.12 arm64} in a \textit{single AWS region} (eu-west-1). The experiment does not extend to:

\begin{itemize}
    \item Multi-item transactions (DynamoDB \texttt{TransactWriteItems}), which incur higher capacity costs and exhibit different failure modes \citep{Idziorek2023DynamoDBTx}.
    \item Provisioned capacity mode, which may exhibit different latency characteristics under throttling.
    \item Alternative managed platforms (Google Cloud Functions, Azure Functions) or alternative NoSQL stores (Amazon S3, Google Cloud Storage).
    \item Cross-region replication or multi-region consistency models.
\end{itemize}

The injected timeout mechanism forces a failure \textit{after} the DynamoDB write has committed, simulating the duplicate-generating case. Real-world failures (partial writes, network partitions mid-transaction) may exhibit different behaviour; the study measures post-commit retry correctness only.

\section{Report Structure}

The remainder of this report is organised as follows. Chapter 2 surveys the literature on serverless fault tolerance, custom runtimes, transactional guarantees, and managed platform primitives, identifying the research niche. Chapter 3 details the experimental methodology, including independent and dependent variables, pilot-sizing rule, statistical tests, and ethics considerations. Chapter 4 describes the design and implementation of the Lambda function, driver, analysis pipeline, and infrastructure-as-code. Chapter 5 presents the evaluation results (currently functional validation on moto; live AWS campaign pending), hypothesis testing, and comparison to the Halfmoon baseline idea. Chapter 6 concludes with a discussion of objectives met, practitioner guidance, and future work.

Vikas Reddy Amanagantti
